\section{Introduction}\label{sec:introduction}

Large language model agents can now sustain mathematical formalization
efforts spanning weeks and more than a hundred pull requests, rather
than single proof-search episodes, shifting the central reliability
question. A proof assistant's kernel guarantees that a \emph{proved}
declaration is correct given its stated type, but kernel checking is
silent on whether that type is the theorem the author actually
intended; whether a previously proved declaration silently disappeared
during a refactor; whether a build genuinely succeeded, versus being
reported successful over a stale disk state; and whether the project's
own development record is itself accurate. These statement-, process-,
and record-level failure modes are the load-bearing risk that remains
once proof-checking is presumed reliable.

We report on a formalization campaign that ran under exactly this
regime. Over roughly six weeks (2026-06-10 through the 2026-07-20
release day), the project consumed, on evidence collected so far, at
least 15.1M output tokens, an API-equivalent cost of \$6{,}972.83, and
199.1\,h of active engagement time under a five-minute inter-event gap
cap -- documented lower bounds, since coverage gaps in the session-log
record remain. % CLAIM: CLM-002
The campaign's central artifact is a kernel-only Lean~4~\cite{moura2021lean4}
formalization, built on mathlib~\cite{mathlib2020}, of existence of
Leray--Hopf weak solutions~\cite{leray1934,hopf1951} to the homogeneous,
incompressible Navier--Stokes equations, comprising two capstone existence
theorems -- one on the periodic 3-torus $\mathbb{T}^3$, one on
$\mathbb{R}^3$ -- whose \texttt{\#print axioms} output is exactly the
three standard Lean kernel axioms, with zero project-declared axioms and
no \texttt{sorryAx} on either theorem~\cite{lerayhopf2026repo}.
% CLAIM: CLM-001
To our knowledge, based on a literature and repository survey conducted in
July 2026 and summarized in the related-work discussion below, no
machine-checked existence theory for a nonlinear fluid PDE has previously
been reported in Lean, Coq, or Isabelle/HOL; we state this as a bound on
what our survey found, not as a proof of non-existence.\footnote{Search
terms, dates, and the nearest adjacent results are recorded in the
project's related-work notes and re-checked immediately before
submission.}
% CLAIM: CLM-006

Three literatures bound this work. Among AI-authored Lean projects,
TauCeti~\cite{tauceti2026}, incubated by the Lean FRO and the Mathlib
Initiative, pursues breadth over a thirteen-area roadmap under human
governance, but reports no accompanying paper, incident-level account,
or session-level evidence trail; we instead pursue depth against a
single frozen scope, with a named-incident case study as a central
artifact. Among formalized analysis adjacent to fluid PDE, mathlib's
Gagliardo--Nirenberg--Sobolev inequality~\cite{vandoorn2024gns} and a
De Giorgi--Nash--Moser elliptic regularity
formalization~\cite{armstrong2026degiorginashmoser} establish a
building-block inequality and a regularity theory respectively, neither
an existence theorem for a nonlinear PDE; the nearest precedent, an
AI-assisted formalization of well-posedness for the linear, kinetic
Vlasov equation~\cite{miller2026vlasov}, is single-author and
single-agent, with no multi-agent governance or incident structure.
Large-scale human-team formalizations such as the Odd Order
Theorem~\cite{gonthier2013oddorder} and the Liquid Tensor
Experiment~\cite{commelin2023abstraction} manage complexity through a
human team and a specification-driven blueprint; our setting substitutes
a single scientific owner governing agent-authored implementation, with
a frozen scope document as blueprint. Finally, the agentic and LLM
theorem-proving
literature~\cite{xin2025deepseekproverv2,yang2023leandojo,deepmind2025alphaproof,zhou2026leanatlas}
reports session- or benchmark-scale results, not a project-scale
process with recorded incidents; rising autonomous task
horizons~\cite{metr2025longtasks} and a study of agentic pull
requests~\cite{agenticprs2026security} motivate the reliability question
from outside formalization but do not instrument a single campaign to
this depth.

This paper makes three contributions:
\begin{enumerate}
\item A completed, kernel-only Lean~4 and mathlib formalization of
  Leray--Hopf weak-solution existence on $\mathbb{T}^3$ and $\mathbb{R}^3$,
  under an agent-orchestrated regime over roughly six weeks.
\item A workflow design -- statement-first development, independent
  adversarial review, and machine-checked verification gates -- with an
  evidence-backed account of what it caught and missed, drawn from named
  incidents rather than impressions.
\item An evidence methodology for AI-assisted research: session-log
  escrow with a manifest, machine-checked claim-to-evidence verification,
  and cost accounting reconciled against billing records.
\end{enumerate}

Section~\ref{sec:formalization} states the two capstone theorems.
Section~\ref{sec:workflow} describes the agent-orchestrated workflow.
Section~\ref{sec:incidents} presents four incidents where a build or
axiom check alone was insufficient to guarantee soundness, and what
closed the gap. Section~\ref{sec:discussion} discusses the resulting
governance picture and limitations; Section~\ref{sec:conclusion}
concludes.
